\section{Fixed- and minimum-credit energy games}
\label{sec:energy-credits}

The preceding algorithms also compute the initial credits needed in
one-resource energy games. Retain the game model of
Theorem~\ref{thm:main}. For a start $v\in V$ and an integer credit
$c\ge0$, call
$(v,c)$ \emph{energy-winning} if Max has a strategy such that, against
every history-dependent Min strategy, the induced play satisfies
\begin{equation}
 c+\sum_{t=0}^{j-1}w(e_t)\ge0\qquad\text{for every }j\ge0.
 \label{eq:energy-objective}
\end{equation}
There is one accumulated resource and no upper bound or saturation.
Write $c_{\min}(v)$ for the least nonnegative integer credit for which
$(v,c)$ is energy-winning, and put $c_{\min}(v)=\infty$ if there is
none. A positional Max policy specifies one labelled outgoing edge at
each Max vertex.

The reductions connecting these problems to mean-payoff games are
classical. Bouyer, Fahrenberg, Larsen, Markey, and Srba formulate the
fixed-credit lower-bound problem and reduce it to mean payoff using
adversary-controlled returns to the start
\cite[Section~2.1 and Proposition~12]{BouyerEtAl2008}. Our construction
places the initial credit on each return edge. Brim, Chaloupka, Doyen,
Gentilini, and Raskin distinguish the unknown-credit and minimum-credit
problems, record the unknown-credit/mean-payoff equivalence, and
characterize minimum credits by an energy progress measure
\cite[Section~2 and Theorems~3 and~5]{BrimEtAl2011}. The reset construction
below lets us apply Corollary~\ref{cor:full-solution} to a prescribed-credit
instance, recovering a policy when that instance is winning. A bound on
finite minimum credits then permits exact binary search.

\begin{corollary}[Fixed and minimum energy credits]
\label{cor:energy-credits}
Uniform deterministic algorithms have the following guarantees.
\begin{enumerate}
\item Given $v$ and a nonnegative binary integer $c$, decide whether
$(v,c)$ is energy-winning. On a positive answer, return a positional
Max policy satisfying Equation~\eqref{eq:energy-objective} from that
start and credit against every history-dependent Min strategy.
\item Given $v$, return the exact $c_{\min}(v)$, including the possible
answer $\infty$. On a finite answer, return a positional Max policy
winning from $v$ at credit $c_{\min}(v)$.
\end{enumerate}
The fixed-credit computation takes
$2^{O((\log_2(L_c+2))^2)}$ bit operations, where $L_c$ is the full
explicit input length including the start and credit. The
minimum-credit computation takes $2^{O((\log_2(L+2))^2)}$ bit
operations for the original game length $L$. Handling every vertex by
separate per-start searches, including a separate policy
for every finite answer, has the same form of total bit bound.

Each policy may depend on the queried start and credit; it also wins
from that same start at every larger credit.
\end{corollary}

\begin{proof}
\noindent\textbf{Fixed credit by resetting.}\enspace
For the fixed pair $(v,c)$ construct a mean-payoff game $R_{v,c}$.
Replace every $u\in V$ by an entry $u_{\mathrm{in}}$ owned by Min and
an exit $u_{\mathrm{out}}$ owned by the original owner of $u$. Add the
two edges
\[
 u_{\mathrm{in}}\xrightarrow{c}v_{\mathrm{in}},
 \qquad
 u_{\mathrm{in}}\xrightarrow{0}u_{\mathrm{out}},
\]
and replace each labelled original edge $e:u\to z$ by its own edge
\[
 u_{\mathrm{out}}\xrightarrow{w(e)}z_{\mathrm{in}}.
\]
Start at $v_{\mathrm{in}}$. The game is total and has $2n$ vertices
and $m+2n$ edges, where $m=|E|$. The reset at $u=v$ is a nonnegative
self-loop, including when $c=0$, and distinct original parallel edges
remain distinct. If $U(R_{v,c})$ denotes its zero-threshold winning
set, then
\begin{equation}
 (v,c)\text{ is energy-winning in }G
 \quad\Longleftrightarrow\quad v_{\mathrm{in}}\in U(R_{v,c}).
 \label{eq:energy-reset-equivalence}
\end{equation}

For the forward implication, use an energy-winning strategy from
$(v,c)$ and restart its history after each reset. In any transformed
prefix, every completed segment has original-edge sum at least $-c$
and then receives the reset weight $+c$. Its total is therefore
nonnegative. The unfinished segment has sum at least $-c$; the
intervening zero edges do not change this bound. Every transformed
prefix sum is consequently at least $-c$. Dividing by its length
shows that the lower limiting average is nonnegative. This argument
also covers no resets and infinitely many resets, and does not assume
that the original energy strategy is positional.

Conversely, if $v_{\mathrm{in}}\in U(R_{v,c})$,
Corollary~\ref{cor:full-solution} returns a positional Max strategy
for $R_{v,c}$ guaranteeing nonnegative mean payoff from that start.
Restrict its choices at Max exits to obtain a positional policy in
$G$. Suppose a consistent original prefix of length $j$ from $v$ had
weight $q<-c$. In the reset game Min can choose continuation at its
entries, follow that prefix at its owned exits, and then reset. This
is a closed walk from $v_{\mathrm{in}}$ with weight $q+c<0$ and
length $2j+1$. Positionality of the Max strategy makes the same walk
consistent on every repetition. Min could therefore repeat it forever,
giving mean payoff $(q+c)/(2j+1)<0$, a contradiction. Thus all original
prefixes have weight at least $-c$. This proves the equivalence and
the positive-answer policy guarantee; only the Max strategy from the
full-solution call is used.

\par\medskip\noindent\textbf{A finite credit bound and exact search.}\enspace
For the original weights put
\begin{equation}
 W_{\mathrm{in}}=\max_{e\in E}|w(e)|,
 \qquad H=(n-1)W_{\mathrm{in}}.
 \label{eq:energy-credit-bound}
\end{equation}
In particular $W_{\mathrm{in}}$ and $H$ may be zero. An energy winner
with any finite credit is an original zero-threshold mean-payoff
winner, since its prefix sums are bounded below. Conversely, the
cycle-deletion estimate in the proof of
Proposition~\ref{prop:game-reduction} gives a positional Max policy
on the original winning region for which every consistent prefix has
weight at least $-H$. Indeed, its permitted simple cycles have
nonnegative original sum, and deleting such cycles from a prefix leaves
at most $n-1$ edges, each of weight at least $-W_{\mathrm{in}}$.
Consequently, for the original winning set $U(G)$,
\begin{equation}
 c_{\min}(v)<\infty\quad\Longleftrightarrow\quad v\in U(G),
 \qquad
 v\in U(G)\quad\Longrightarrow\quad 0\le c_{\min}(v)\le H.
 \label{eq:energy-finite-equivalence}
\end{equation}
The algorithm uses this existence bound as a known winning endpoint;
it need not compute the potential used in that proposition.

Compute $U(G)$ once by Theorem~\ref{thm:main} and output $\infty$ for
starts outside it. For a winning start, the fixed-credit predicate is
monotone: a strategy that wins at $c$ still wins at every larger
credit. Binary search with a formal losing lower endpoint $\ell=-1$
and the known winning upper endpoint $h=H$. While $h-\ell>1$, query
$\lfloor(\ell+h)/2\rfloor$ by
Equation~\eqref{eq:energy-reset-equivalence} and replace the endpoint
having the same answer. The sentinel itself is never queried. This takes at most
$\lceil\log_2(H+1)\rceil$ fixed-credit decisions and returns the exact
least winning integer. When $H=0$ the initial gap is one, so no search
query is made and the answer is zero. This includes both $n=1$ and
all-zero weights; any start without a finite credit has already been
excluded by $U(G)$. To return a policy at the minimum, apply
Corollary~\ref{cor:full-solution} to $R_{v,c_{\min}(v)}$ and restrict
its Max choices as above.

\par\medskip\noindent\textbf{Bit complexity and output size.}\enspace
Let $b_c$ be the binary length of the supplied nonnegative credit,
so $b_c=O(1+\log_2(c+1))$. The explicit reset game has encoding length
\[
 O\bigl(L+n(b_c+\log_2(n+1))\bigr),
\]
which is polynomial in $L_c$. The credit is copied into only $n$ edge
records; no array indexed by its numerical value is constructed.
Theorem~\ref{thm:main} decides membership in its winning set, and a
single application of Corollary~\ref{cor:full-solution} on a positive
answer supplies the policy, within the fixed-credit bound claimed.

For minimum credits, $H$ and every queried credit have
$O(1+\log_2(n+1)+\log_2(W_{\mathrm{in}}+1))=O(L+1)$ bits. All queried
encodings therefore have polynomial length in $L+1$. For all vertices,
the outer reduction makes at most $1+n\lceil\log_2(H+1)\rceil$ direct
winning-set calls, including the original one, and at most one
full-solution call for each requested finite-credit policy. These
counts are polynomial in
$L+1$. Polynomially many calls on polynomial-length inputs preserve
$2^{O((\log_2(L+2))^2)}$ bit complexity. Computing $H$, the binary
search arithmetic, and copying the policies have polynomial bit cost.
The output also has polynomial length: each finite credit has the
stated bit bound, and each per-start policy lists one original edge
at each Max vertex. No energy-indexed array is used in the search.
\end{proof}
